\SourceLeaf{244}%
gegebenen Fläche liegen soll, deren Beschaffenheit durch die Gleichung \(P\mathrm{d}x+Q\mathrm{d}y+R\mathrm{d}z=0\) ausgedrückt wird, müssen auch die Variationen \(\delta x,\delta y,\delta z\) der Gleichung \(P\delta x+Q\delta y+R\delta z=0\) genügen. Aus bekannten Prinzipien folgt leicht, dass die Differentiale
\[
\mathrm{d}\frac{\mathrm{d}x}{\BareRoot{(\mathrm{d}x^2+\mathrm{d}y^2+\mathrm{d}z^2)}},\quad
\mathrm{d}\frac{\mathrm{d}y}{\BareRoot{(\mathrm{d}x^2+\mathrm{d}y^2+\mathrm{d}z^2)}},\quad
\mathrm{d}\frac{\mathrm{d}z}{\BareRoot{(\mathrm{d}x^2+\mathrm{d}y^2+\mathrm{d}z^2)}}
\]
der Reihe nach zu den Größen \(P,Q,R\) proportional sein müssen. Nun sei \(\mathrm{d}r\) das Element der gekrümmten Linie, \(\lambda\) der Punkt auf der Kugeloberfläche, der die Richtung dieses Elements darstellt, und \(L\) der Punkt auf der Kugeloberfläche, der die Richtung der Normalen zur gekrümmten Fläche darstellt. Schließlich seien \(\xi,\eta,\zeta\) die Koordinaten des Punktes \(\lambda\) und \(X,Y,Z\) die Koordinaten des Punktes \(L\), jeweils bezogen auf den Mittelpunkt der Kugel. Damit gilt
\[
\mathrm{d}x=\xi\mathrm{d}r,\quad\mathrm{d}y=\eta\mathrm{d}r,\quad\mathrm{d}z=\zeta\mathrm{d}r
\]
Daraus schließen wir, dass sich für jene Differentiale \(\mathrm{d}\xi,\mathrm{d}\eta,\mathrm{d}\zeta\) ergeben. Da die Größen \(P\mathord{.}\ Q\mathord{.}\ R\) zu \(X,Y,Z\) proportional sind, wird die kürzeste Linie durch die Gleichungen gekennzeichnet
\[
\frac{\mathrm{d}\xi}{X}=\frac{\mathrm{d}\eta}{Y}=\frac{\mathrm{d}\zeta}{Z}
\]
Ferner ist leicht einzusehen, dass \(\BareRoot{(\mathrm{d}\xi^2+\mathrm{d}\eta^2+\mathrm{d}\zeta^2)}\) dem kleinen Bogen auf der Kugeloberfläche gleich ist, der den Winkel zwischen den Richtungen der Tangenten am Anfang und am Ende des Elements \(\mathrm{d}r\) misst. Der Ausdruck ist also \(=\frac{\mathrm{d}r}{\rho}\), wenn \(\rho\) den Krümmungsradius der kürzesten Linie an dieser Stelle bezeichnet. Somit gilt
\[
\rho\mathrm{d}\xi=X\mathrm{d}r,\quad\rho\mathrm{d}\eta=Y\mathrm{d}r,\quad\rho\mathrm{d}\zeta=Z\mathrm{d}r
\]
\Article{15}
Nehmen wir an, dass auf einer gekrümmten Fläche von einem gegebenen Punkt \(A\) unzählige kürzeste Kurven ausgehen. Wir unterscheiden sie durch den Winkel, den das Anfangselement jeder dieser Kurven mit dem Anfangselement einer als erste gewählten Linie bildet. Es sei \(\varphi\) dieser Winkel oder allgemeiner eine Funktion davon, und \(r\) sei die Länge einer solchen kürzesten Linie vom Punkt \(A\) bis zum Punkt mit den Koordinaten \(x,y,z\). Da bestimmten Werten der Veränderlichen \(r,\varphi\) bestimmte Punkte der Fläche entsprechen, lassen sich die Koordinaten \(x,y,z\) als Funktionen von \(r,\varphi\) auffassen. Wir behalten die Bezeichnungen \(\lambda,L,\xi,\eta,\zeta,X,Y,Z\) in derselben Bedeutung bei
